% Acronyms
%\newacronym{api}{API}{Application Program Interface}

% Glossary
% \newglossaryentry{apig}{
%     name={API},
%     text={Application Program Interface},
%     sort=api,
%     description={In informatics, an API is a set of procedures available to programmers, typically grouped to form a toolkit for a specific task within a program. Its purpose is to provide an abstraction, usually between hardware and the programmer or between low-level and high-level software, simplifying the programming process}
% }

%####################################################################################
% GLOSSARIO NORMALE
%####################################################################################
\newglossaryentry{embedded_g}{
    name={Embedded},
    text={embedded},
    sort={Embedded},
    description={Indica un sistema integrato \textit{hardware-software} progettato per svolgere una funzione specifica e dedicata all'interno di un dispositivo. I sistemi \textit{embedded} operano tipicamente con risorse limitate (CPU, memoria, energia) e in ambiti come: centraline auto, elettrodomestici, dispositivi IoT, \textit{wearable} e strumentazione medica}
}
\newglossaryentry{func_saf_g}{
    name={Functional Safety},
    text={functional safety},
    sort={Functional Safety},
    description={Ramo della sicurezza che si occupa di garantire il corretto comportamento di un sistema o uno strumento in seguito al cambiamento dei parametri di \textit{input}, o al verificarsi di errori e/o guasti \textit{hardware}}
}
\newglossaryentry{agile_g}{
    name={Agile},
    text={agile},
    sort={Agile},
    description={Fondato sui principi dell'\textit{Agile Manifesto} (2001), promuove il coinvolgimento continuo del cliente, il rilascio frequente di valore funzionante e l'auto-organizzazione dei \textit{team}. \textit{Framework} come \textit{Scrum} e \textit{Kanban} ne rappresentano applicazioni pratiche diffuse}
}
\newglossaryentry{scrum_g}{
    name={Scrum},
    text={scrum},
    sort={Scrum},
    description={\textit{Framework} che applica i principi \textit{Agile} per la gestione e lo sviluppo di prodotti complessi, organizzato in cicli temporali brevi e a durata fissa chiamati \textit{Sprint} (tipicamente 2-4 settimane).
    Una figura chiave è lo \textit{scrum master}, responsabile di facilitare l'adozione e la corretta applicazione del \textit{framework Scrum} all'interno del \textit{team}.
    Egli rimuove ostacoli che rallentano il lavoro, facilita le cerimonie \textit{Scrum} (\textit{Daily, Sprint Planning, Review, Retrospective}) e isola il \textit{team} da interferenze esterne, promuovendo il miglioramento continuo dei processi}
}
\newglossaryentry{open_source_g}{
    name={Open Source},
    text={open source},
    sort={Open Source},
    description={Prodotto \textit{software} in cui il codice sorgente è reso pubblicamente disponibile, consentendo a chiunque di visualizzarlo, modificarlo e ridistribuirlo}
}
\newglossaryentry{ml_g}{
    name={Machine Learning (ML)},
    text={machine learning},
    sort={Machine learning},
    description={Branca dell'intelligenza artificiale che studia algoritmi capaci di apprendere \textit{pattern} dai dati, migliorando le proprie prestazioni su un compito senza essere esplicitamente programmati per eseguirlo}
}
\newglossaryentry{isotp_g}{
    name={ISO-TP / ISO 15765-2},
    text={IsoTp},
    sort={IsoTp},
    description={Protocollo di livello trasporto definito dallo \textit{standard} ISO 15765-2, usato principalmente nel settore \textit{automotive} per trasmettere messaggi su \textit{bus} CAN.
    Poiché i \textit{frame} CAN \textit{standard} hanno un \textit{payload} limitato a 8 \textit{byte}, ISO-TP permette di segmentare messaggi più grandi in più \textit{frame} e di riassemblarli a destinazione, gestendo anche il controllo di flusso tra trasmettitore e ricevitore}
}
\newglossaryentry{modbus_g}{
    name={Modbus},
    text={Modbus},
    sort={Modbus},
    description={Protocollo di comunicazione sviluppato originariamente per l'automazione industriale, basato su un'architettura \textit{master/slave}.
    Permette a un dispositivo \textit{master} di leggere e scrivere dati (registri, \textit{coil}, ingressi digitali/analogici) su uno o più dispositivi \textit{slave} tramite richieste e risposte strutturate in un formato semplice e aperto}
}
%####################################################################################
% ELENCO ACRONIMI
%####################################################################################
\newacronym{ReD}{R\&D}{Research and Development}
\newacronym{po}{PO}{Product Owner}
\newacronym{OEM}{OEM}{Original Equipment Manufacturer}
\newacronym{misra}{MISRA}{Motor Industry Software Reliability Association}
\glsunset{misra}
\newacronym{IDE}{IDE}{Integrated Development Environment}
\glsunset{IDE}
\newacronym{ci_cd}{CI/CD}{Continuous Integration/Continuous Delivery}
\glsunset{ci_cd}
\newacronym{cli}{CLI}{Command Line Interface}
\glsunset{cli}
\newacronym{poc}{PoC}{Proof of Concept}
\newacronym{cra}{CRA}{Cyber Resiliance Act}
\newacronym{csms}{CSMS}{Cyber Security Management System}
\newacronym{IoT}{IoT}{Internet of Things}
\glsunset{IoT}
\newacronym{can}{CAN}{Controller Area Network}
\glsunset{can}
\newacronym{api}{API}{Application Programming Interface}
\glsunset{api}
\newacronym{uds}{UDS}{Unified Diagnostic Services}
\newacronym{uart}{UART}{Universal Asynchronous Receiver-Transmitter}
\glsunset{uart}

%####################################################################################
% GLOSSARIO ESTESO DEGLI ACRONIMI
%####################################################################################

\newglossaryentry{po_ext}{
    name={Product Owner},
    text={PO},
    sort={Product Owner},
    description={Figura chiave nei metodo \textit{Agile}, responsabile di massimizzare il valore del prodotto sviluppato dal \textit{team}.
    Definisce la visione del prodotto, gestisce e prioritizza le attività, e traduce le esigenze di \textit{stakeholder} e clienti in requisiti chiari per il \textit{team}. I suoi compiti possono variare in base all'organizzazione della realtà in cui si trova, andando a sovrapporsi in parte con quelli del \textit{Project Manager}}
}
\newglossaryentry{misra_ext}{
    name={MISRA},
    text={MISRA},
    sort={MISRA},
    description={Organizzazione che pubblica linee guida di programmazione (principalmente per C e C++) volte a migliorare sicurezza, affidabilità e portabilità del codice, soprattutto in ambito \textit{automotive} ed \textit{embedded/safety-critical}.
    Gli \textit{standard} MISRA definiscono regole e \textit{best practice} per evitare costrutti ambigui o pericolosi del linguaggio, e sono spesso richiesti in contesti normativi o di certificazione}
}
\newglossaryentry{IDE_ext}{
    name={IDE},
    text={IDE},
    sort={IDE},
    description={Ambiente di Sviluppo Integrato, \textit{software} che riunisce in un'unica interfaccia gli strumenti necessari alla scrittura di codice}
}
\newglossaryentry{ci_cd_ext}{
    name={CI/CD (Continuous Integration/Continuous Delivery)},
    text={CI/CD},
    sort={CI/CD},
    description={Pratica \textit{DevOps} che automatizza le fasi di integrazione, \textit{test} e rilascio del \textit{software}.
    La \textit{Continuous Integration} prevede l'unione frequente del codice in un \textit{repository} condiviso con \textit{build} e \textit{test} automatici, per individuare errori precocemente.
    La \textit{Continuous Delivery/Deployment} estende il processo automatizzando il rilascio in produzione, riducendo tempi e rischi legati alle \textit{release} manuali}
}
\newglossaryentry{PoC_ext}{
    name={PoC},
    text={PoC},
    sort={PoC},
    description={\textit{Proof of Concept}, si tratta di un prodotto che ha lo scopo di dimostrare la fattibilità di un progetto andando ad implementare, e mettendo in relazione, diverse tecnologie e idee sulle quali esso si basa.
    Può successivamente divenatare una base iniziale sulla quale andare a sviluppare il progetto vero e proprio}
}
\newglossaryentry{cra_ext}{
    name={CRA (Cyber Resilience Act)},
    text={CRA},
    sort={CRA},
    description={Regolamento dell'Unione Europea che introduce requisiti obbligatori di cybersicurezza per i prodotti con elementi digitali (\textit{hardware} e \textit{software}) immessi sul mercato europeo.
    Impone ai produttori obblighi come la gestione delle vulnerabilità, l'aggiornamento sicuro del \textit{software} e la fornitura di informazioni chiare agli utenti sulla sicurezza del prodotto, con l'obiettivo di ridurre il numero di dispositivi vulnerabili e responsabilizzare i produttori anche dopo la vendita.}
}
\newglossaryentry{can_ext}{
    name={CAN},
    text={CAN},
    sort={CAN},
    description={Protocollo di comunicazione seriale \textit{multi-master}, utilizzato per la comunicazione tra ECU \textit{(Electronic Control Units)} in ambito \textit{automotive}.
    Le varie ECU sono collegate fisicamente tramite il \textit{CAN bus}, e comunicano mandano messaggi \textit{broadcast} nella rete condivisa}
}
\newglossaryentry{api_ext}{
    name={API},
    text={API},
    sort={API},
    description={Interfaccia che mette a disposizione una serie di servizi tramite cui componenti \textit{software} distinti possono comunicare e scambiarsi dati, senza necessità di conoscere i dettagli implementativi reciproci}
}
\newglossaryentry{uds_ext}{
    name={UDS (Unified Diagnostic Services)},
    text={UDS},
    sort={UDS},
    description={Anche conosciuto come ISO 14229, è un protocollo a livello applicativo usato nel settore \textit{automotive} per la comunicazione diagnostica tra un \textit{tool} esterno e le centraline elettroniche (ECU) di un veicolo.
    Definisce un insieme di servizi identificati da codici esadecimali scambiati tramite richieste e risposte strutturate.
    UDS si appoggia tipicamente a protocolli di trasporto come ISO-TP (ISO 15765-2) per la trasmissione dei messaggi su \textit{bus} CAN}
}
\newglossaryentry{uart_ext}{
    name={UART (Universal Asynchronous Receiver-Transmitter)},
    text={UART},
    sort={UART},
    description={Periferica \textit{hardware} per la comunicazione seriale asincrona punto-punto, ampiamente usata nei sistemi \textit{embedded}.
    Trasmette dati \textit{bit} per \textit{bit} su due linee separate, sincronizzandosi tramite \textit{bit} di \textit{start/stop} e una velocità di trasmissione (\textit{baud rate}) concordata a priori tra i due dispositivi}
}
